\chapter{Meininga som fundament}

\begin{motto}
Michl angrip ikkje Sullivan. Han angrip ein\\ karikatur av Sullivan som modernismen sjølv laga.
\end{motto}

Den hardaste enkeltsetninga i \textsc{Grunnlaust} er at «form følgjer funksjon» er ein tautologi — eit ord som ser ut til å seie noko, men ikkje gjer det, fordi «funksjon» kan tøyast til å tyde kva som helst designaren alt har bestemt seg for. Om huset har eit flatt tak, var funksjonen flatt tak; om det har spiss gavl, var funksjonen spiss gavl. Formelen forklarar då ingenting; han pyntar berre på vala etterpå. Og om dette stemmer, er det alvorleg, for då er sjølve grunnsetninga i moderne design tom, og eit fag som kviler på ein tom setning, kviler ikkje på noko. Dette kapitlet bestrider påstanden frå to kantar. Det viser at funksjonalismen aldri tydde det karikaturen tillegg han, og at design dessutan har eit fundament som ligg djupare enn funksjonen: meininga. Langt frå å vere tomt har faget eit stoff som er reelt nok til å bere ein heil teori — og det stoffet er kva ting \emph{tyder} for menneske.

\section{Les Sullivan på nytt}

Vidar Michl si kjende avlivinga av «form følgjer funksjon» råkar eit verkeleg mål, men ikkje det målet han trur. Han råkar den modernistiske parolen — den utvatna, dogmatiske versjonen som funksjonalistane på 1920- og 30-talet gjorde til ein regel: pynt er brotsverk, alt overflødig skal bort, forma skal vere den nakne avleiinga av nytta. Den versjonen er verkeleg sårbar for tautologi-skuldinga, og Michl har rett mot han. Men parolen var ikkje Sullivan sin. Går ein til kjelda, til Louis Sullivan sin essay frå 1896, finn ein noko ganske anna \autocite{sullivan1896tall}.

Sullivan skreiv ikkje om nytte. Han skreiv om \emph{uttrykk}. Setninga hans var ikkje ein nøktern ingeniørregel, men ei nesten religiøs påstand om livet sjølv: at i all organisk natur tek forma som ein ting har, etter funksjonen — livsfunksjonen, vesenet — til den tingen. Ørna har si form fordi ho er ein ørn; eikja si fordi ho er ei eik; ein elv, eit skydrag, eit menneske, kvart av dei ber ei form som er det ytre uttrykket for det indre livet. Sullivan kalla det «the law of all organic and inorganic things», og då han bruka det på bygningen, meinte han at ein høg kontorbygning skulle uttrykkje si eiga sjel, sin eigen karakter som høg ting, gjennom forma si. «Funksjon» tyder her ikkje «kva han skal brukast til», men «kva han \emph{er}» — vesenet hans, det han uttrykkjer. Sullivan var ikkje ingeniør i denne setninga; han var romantikar, nær Whitman og den amerikanske transcendentalismen.

Då fell tautologi-skuldinga mot originalen, for ho føreset at «funksjon» tyder nytte, og at nytta så vert tøygd til å passe forma. Men Sullivan sin funksjon var ikkje nytte; han var uttrykk — og uttrykk kan ein lykkast eller mislykkast med. Ein bygning kan uttrykkje vesenet sitt klart eller uklart, ærleg eller falskt, og det er ein dom som ikkje er tom. Michl angrip ikkje Sullivan. Han angrip ein karikatur av Sullivan som modernismen sjølv laga — den same operasjonen \textsc{Grunnlaust} gjentek når boka tek den modernistiske parolen for å vere designfaget sitt grunnlag. Les ein originalen frå 1896, og funksjonen er ikkje nytte; han er uttrykk. Og eit fag som handlar om uttrykk, om kva ting seier og korleis dei seier det, er ikkje eit fag utan innhald.

\section{Funksjonalismen hadde alltid eit innhald}

Sjølv om ein held seg til funksjon i den nøkterne tydinga, fell ikkje skuldinga om tomheit. To kjelder viser kvifor. Den fyrste er Edward De Zurko, som alt i 1957 synte at funksjonalismen ikkje var ein modernistisk oppfinning, men ein tankefigur med ein lang og rik pedigree — tilbake til antikken, gjennom Vitruvius og dei kristne tenkjarane, gjennom opplysningstida \autocite{dezurko1957origins}. I to tusen år har menneske meint at det er ein samanheng mellom kva ein ting er til, og kva form han bør ha, og dei har meint det på mange ulike måtar, med ulikt filosofisk grunnlag. Ein tanke som har vore tenkt så mange gonger, av så ulike hovud, er ikkje ein tautologi. Tautologiar har inga historie, for dei seier ikkje noko ein kan vere usamd om. Funksjonalismen har ei lang historie full av usemje — og det er sjølve prov på at han hevdar noko.

Den andre kjelda gjer poenget skarpt. Glenn Parsons og Allen Carlson, i \emph{Functional Beauty}, tek opp det filosofiske spørsmålet om funksjon og skjønnheit verkeleg heng saman, eller om koplinga berre er ein talemåte \autocite{parsons2008functional}. Svaret deira er at koplinga er ekte. Når vi ser ein ting som er godt eigna til det han er, opplever vi ofte ei eiga form for skjønnheit — ikkje trass i funksjonen, men gjennom han. Eit verktøy som ligg rett i handa, eit fly hvis form fortel om lufta det skjer gjennom, eit dyr formfullenda til sitt levesett: i alle desse ser vi at det å «sjå funksjonell ut» er ein estetisk kvalitet i seg sjølv. Men — og dette er det avgjerande mot tautologi-skuldinga — denne skjønnheita kan ein \emph{ta feil om}. Ein ting kan sjå funksjonell ut utan å vere det, eller vere funksjonell utan å sjå det. Det finst ein skilnad mellom verkeleg og tilsynelatande funksjonell skjønnheit, og å sjå skilnaden krev kunnskap om kva tingen faktisk er til. Då er funksjon-form-koplinga ikkje ein sirkel der alt passar med alt; ho er ein påstand som kan vere sann eller falsk, treffande eller bomma. Og ein påstand som kan vere falsk, er ikkje tom.

Dette er meir enn ein filosofisk finesse. Det gjev designaren sin dom eit objekt. Når ein designar seier at ei form er rett for funksjonen, seier han ikkje «denne forma er den eg valde»; han seier «denne forma uttrykkjer og tener det tingen er til, betre enn alternativa» — og det er noko ein kan grunngje, prøve, og vere usamd om. Tautologien finst berre i karikaturen.

\section{Krippendorff: meining som nytt fundament}

Men det djupaste svaret er ikkje å forsvare funksjonen; det er å vise at design kviler på noko endå meir grunnleggjande. Klaus Krippendorff gjorde dette eksplisitt i \emph{The Semantic Turn} — og undertittelen er ikkje til å misforstå: \emph{A New Foundation for Design} \autocite{krippendorff2006}. Krippendorff sitt utgangspunkt er ei setning som har blitt eit slagord, men som han meinte heilt bokstaveleg: design er å \emph{gjere ting meiningsfulle}, å skape meining — «making sense of things». Det er ikkje funksjonen som er det fyrste; det er meininga. Eit menneske møter aldri ein ting som rein funksjon. Han møter han som noko som tyder noko — som vekkjer ei forventning, ber ei historie, høyrer til ein samanheng, fortel kven brukaren er. Ei dør ein \emph{ser} kan opnast, ein knapp ein \emph{forstår} kva gjer, eit produkt som \emph{seier} kva det er for: alt dette er meining, ikkje funksjon i snever forstand, og det er dette designaren faktisk arbeider med.

Krippendorff sitt poeng er at meining ikkje er pynt lagt oppå funksjonen, men det laget mennesket faktisk lever i. Vi handlar ikkje ut frå kva ting er; vi handlar ut frå kva dei tyder for oss. Difor er det å forme meining — å gjere ting forståelege, brukbare, vedkomande — ikkje ein tilleggsdisiplin, men sjølve designaren si oppgåve. Og dette er eit \emph{fundament} i den fullaste tyding: eit eige objekt (meininga ting ber), ein eigen aktivitet (å gjere ting meiningsfulle), og eit eige kriterium (om meininga lèt seg gripe og leve med). Tanken var ikkje ny i 2006; Krippendorff hadde lagt grunnen alt i arbeidet om produktsemantikk saman med Reinhart Butter, der dei utforska korleis forma til eit produkt ber symbolske kvalitetar — korleis ting kommuniserer kva dei er og kva dei vil med oss \autocite{krippendorff1984semantics}. Det semantiske var med andre ord ikkje ein vri, men ei mogning av ei innsikt faget hadde arbeidd med i tiår.

Det avgjerande mot tomheit-skuldinga er at meininga gjev design eit innhald funksjonen aleine ikkje kan gje. To produkt kan ha nett same funksjon og likevel tyde heilt ulike ting — det eine trygt, det andre skremmande; det eine for fagfolk, det andre for born; det eine som høyrer heime, det andre som er framand. Heile denne skilnaden er usynleg for den som berre ser funksjon, og han er nett der designaren sitt arbeid ligg. Eit fag som har meininga til sitt objekt, står ikkje i eit tomrom. Det står midt i det laget av røynda der menneske faktisk lever.

\section{Verganti: meininga som innovasjon}

At meining er eit reelt fundament og ikkje ein talemåte, syner seg klårast der det får følgjer ein kan måle: i økonomien. Roberto Verganti har vist at dei djupaste innovasjonane ofte ikkje er endringar i kva eit produkt \emph{gjer}, men i kva det \emph{tyder} \autocite{verganti2009}. Ei lampe som forandrar seg frå reiskap til stemningsskapar, eit spel som blir frå underhaldning til rørsle, ei matvare som går frå metting til oppleving: i alle desse er teknikken ofte den same eller enklare, men meininga er ny — og det er meiningsendringa som skaper verdien. Verganti kallar dette \emph{meiningsdriven innovasjon}, og poenget hans for vår diskusjon er at dette er noko design kan gjere som verken teknologien eller marknadsanalysen kan. Marknadsanalysen spør kva folk alt vil ha, og kan difor berre forsterke meininga som finst; teknologien kan endre kva som er mogleg, men ikkje kva det skal tyde. Å foreslå ein \emph{ny} meining — ein ny måte å forstå kva ein ting er for og kvifor han har verdi — er ein eigen kompetanse, og det er designaren sin.

Verganti gjorde dette skarpare i arbeidet saman med Donald Norman \autocite{norman2014incremental}. Dei skilde mellom to slags endring som lett vert blanda saman: endring i teknologi og endring i meining. Vanleg designforsking, hevda dei, er flink til det inkrementelle — å betre eit produkt steg for steg gjennom å studere brukaren. Men dei radikale skifta, dei som forandrar kva ein ting \emph{er} for menneske, kjem sjeldan frå å spørje brukaren, for brukaren kan ikkje be om ei meining han enno ikkje kjenner. Dei kjem frå ei tolking — frå nokon som ser kva ein ting kunne tyde, og vågar å foreslå det. Dette er ein viktig nyanse, for det viser at meiningsarbeidet ikkje er noko vagt og udisiplinert, men har sin eigen struktur: ein veit no skilnaden mellom å betre det som er, og å endre kva noko tyder, og ein veit at dei to krev ulike framgangsmåtar. Det er ikkje slik tom prosess ser ut. Det er slik eit fag ser ut når det har kartlagt sitt eige stoff.

Då er ringen slutta. Skuldinga var at design kviler på ein tom setning. Men «form følgjer funksjon» tydde aldri det karikaturen sa; funksjonalismen har to tusen års innhald og ei genuin kopling mellom funksjon og skjønnheit; og under funksjonen ligg eit endå djupare fundament — meininga — som har eit eige objekt, ein eigen teori frå Krippendorff, og målbare følgjer hos Verganti. Design er ikkje eit fag utan innhald som pyntar vala sine med ein tom formel etterpå. Det er eit fag hvis stoff er noko av det mest verkelege som finst i menneskelivet: kva ting tyder for oss, og korleis dei kjem til å tyde det.

\vignett

Men meining er alltid meining \emph{for} nokon. Bak både funksjonen og meininga står mennesket dei er til for — og det er mennesket faget til sist svarar overfor. Det er emnet for neste kapittel.

\vidarelesing{\fullcite{krippendorff2006}; \fullcite{verganti2009}; \fullcite{norman2014incremental}; \fullcite{parsons2008functional}; \fullcite{sullivan1896tall}; \fullcite{dezurko1957origins}; \fullcite{krippendorff1984semantics}.}
